\begin{theorem}[The Euclidean additive group on $\mathbb R^N$ is a Carnot setting]
For every $N\in\mathbb N$ there exists a Carnot setting $\mathbb G$ (\noderef{CarnotSetting})
whose data $(n,m,r)$, group law $*$, first-layer projection $\pi_1$, horizontal exponential
$\exp_{\mathrm h}$, homogeneous norm $\|\cdot\|$ and horizontal frame
$\sigma=[\sigma^1,\ldots,\sigma^m]$ satisfy
\begin{enumerate}
\item $n=N$, $m=N$ and $r=1$;
\item $x*y=x+y$ for all $x,y\in\mathbb R^{n}$;
\item $\|x\|=|x|$ for all $x\in\mathbb R^{n}$, where $|\cdot|$ is the Euclidean norm;
\item $\exp_{\mathrm h}(\pi_1(x))=x$ for all $x\in\mathbb R^{n}$;
\item $\sigma^i(x)=\exp_{\mathrm h}(e_i)$ for all $x\in\mathbb R^{n}$ and all
  $i\in\{1,\ldots,m\}$, where $e_1,\ldots,e_m$ is the standard basis of $\mathbb R^{m}$.
\end{enumerate}
\end{theorem}
\begin{proof}
Fix $N\in\mathbb N$ and write $e_1,\ldots,e_N$ for the standard basis of $\mathbb R^N$. We
exhibit the data and verify, one at a time, every condition of \noderef{CarnotGroup},
\noderef{CarnotNorm} and \noderef{CarnotSetting}, and then conditions (1)--(5).

\medskip\noindent\textbf{The data.} Put
\[
 n:=N,\qquad m:=N,\qquad r:=1,
\]
\[
 x*y:=x+y,\qquad \pi_1:=\mathrm{id}_{\mathbb R^N},\qquad
 \exp_{\mathrm h}:=\mathrm{id}_{\mathbb R^N},\qquad \|x\|:=|x|,\qquad
 \sigma^i(x):=e_i\ \ (i=1,\ldots,N).
\]
Here $\mathrm{id}_{\mathbb R^N}$ is the identity, which is a continuous linear map
$\mathbb R^N\to\mathbb R^N$, as required of $\pi_1$; and $\sigma$ is the constant frame, so
the associated first-order operators are $X_i=e_i\cdot\nabla=\partial_{x_i}$, which is the
frame the paper attaches to the Euclidean group (Section Section~1.2: ``the case $m=n$
corresponds to $r=1$ and $\mathbb G$ being the Euclidean additive group, with
$X_i=\partial_{x_i}$''). Since $r=1$ we have $q:=2\,r!=2$, and the left-invariant distance
of \noderef{CarnotNorm} is
\[
 d(x,y)=\|(-y)*x\|=|(-y)+x|=|x-y| ,
\]
the Euclidean distance.

\medskip\noindent\textbf{Step 1: the conditions of \noderef{CarnotGroup}.}

\emph{(i) $r\ge1$.} Immediate from $r=1$.

\emph{(ii) $(\mathbb R^N,*)$ is a group with neutral element $0$ and inversion
$x\mapsto-x$.} The five required identities are
$(x+y)+z=x+(y+z)$, $x+0=x$, $0+x=x$, $x+(-x)=0$ and $(-x)+x=0$; these are the axioms of the
additive group $\mathbb R^N$.

\emph{(iii) The group law is $C^\infty$.} The map
$\mathbb R^N\times\mathbb R^N\to\mathbb R^N$, $(x,y)\mapsto x+y$, is the sum of the two
coordinate projections. Each projection is continuous and linear, hence $C^\infty$, and a
sum of $C^\infty$ maps is $C^\infty$.

\emph{(iv) $\pi_1(x*y)=\pi_1(x)+\pi_1(y)$.} Both sides equal $x+y$ because $\pi_1$ is the
identity.

\emph{(v) $\exp_{\mathrm h}$ is continuous and $\pi_1(\exp_{\mathrm h}(\zeta))=\zeta$.} The
identity is continuous, and $\pi_1\circ\exp_{\mathrm h}=\mathrm{id}$.

\emph{(vi) $r=1$ implies $*$ commutative.} Addition on $\mathbb R^N$ is commutative.

\medskip\noindent\textbf{Step 2: the conditions of \noderef{CarnotNorm}.}

\emph{(i) $\|x\|\ge0$; $\|x\|=0$ iff $x=0$; $\|-x\|=\|x\|$; $\|\cdot\|$ continuous.} These
are the defining properties of the Euclidean norm $|\cdot|$ together with its continuity.

\emph{(ii) $\|\exp_{\mathrm h}(\zeta)\|=|\zeta|$.} Since $\exp_{\mathrm h}$ is the identity
and $\|\cdot\|=|\cdot|$, both sides are $|\zeta|$.

\emph{(iii) $x\mapsto\|x\|^{q}$ is $C^\infty$.} Here $q=2$, and
$|x|^{2}=\sum_{i=1}^{N}x_i^{2}$ is a polynomial in the coordinates of $x$, hence $C^\infty$.

\emph{(iv) First inequality of (2.17).} Given $\nu>0$, take $c:=1>0$. For all $x,y$
(in particular whenever $\|x\|<\nu$ and $\|y\|<\nu$) the triangle inequality for the
Euclidean norm gives $\|x*y\|=|x+y|\le|x|+|y|=c\,(\|x\|+\|y\|)$.

\emph{(v) Second inequality of (2.17).} Given $\nu>0$, take $c:=1>0$. Since addition
is commutative, $x^{-1}*y*x=((-x)+y)+x=y+((-x)+x)=y$, so
$\|x^{-1}*y*x\|=|y|$. As $r=1$ we have $1/r=1$, and $c\,\|y\|^{1/r}=|y|^{1}=|y|$. Hence the
required inequality holds with equality.

\emph{(vi) H\"older equivalence on bounded sets.} Given $\nu>0$, take $c:=1>0$. For all
$x,y$ we computed $d(x,y)=\|(-y)*x\|=|x-y|$, and with $1/r=1$ the two required inequalities
read $c^{-1}|x-y|=|x-y|\le|x-y|$ and $|x-y|\le c\,|x-y|^{1}=|x-y|$; both hold with
equality, in particular whenever $|x|\le\nu$ and $|y|\le\nu$.

\emph{(vii) Balls are non-empty and connected.} Fix $x\in\mathbb R^N$ and $\rho>0$. By the
computation of $d$,
\[
 B_\rho(x)=\{y:\ \|(-x)*y\|<\rho\}=\{y:\ |y-x|<\rho\},
\]
the Euclidean open ball of radius $\rho$ about $x$. It contains $x$, since $|x-x|=0<\rho$,
so it is non-empty; and it is convex, hence connected.

\emph{(viii) Properness.} Let $C\in\mathbb R$. Then
$\{x:\|x\|\le C\}=\{x:|x|\le C\}$ is the closed Euclidean ball of radius $C$ about the
origin (empty if $C<0$), which is bounded.

\medskip\noindent\textbf{Step 3: two computations for the constant frame.} Both are used
repeatedly below.

\emph{(3a) The horizontal gradient reads off the frame directions.} By \noderef{hgrad},
for every $x$, every $1$-jet $D$ and every $i$,
\[
 \big(\mathrm{hgrad}\ \sigma\ x\ D\big)_i=D\big(\sigma^i(x)\big)=D(e_i).
\]
In particular this value does not depend on $x$.

\emph{(3b) The correction matrix vanishes, so the symmetrized horizontal second derivative
is the matrix of the $2$-jet in the frame directions.} By \noderef{Mfield},
\[
 \mathcal M(x,D)_{i,j}=\tfrac12\Big(D\big(D\sigma^{j}(x)\,\sigma^{i}(x)\big)
   +D\big(D\sigma^{i}(x)\,\sigma^{j}(x)\big)\Big).
\]
Each $\sigma^{k}$ is a constant map, so its Fr\'echet derivative $D\sigma^{k}(x)$ is the
zero linear map; hence $D\sigma^{j}(x)\sigma^{i}(x)=0$ and
$D\sigma^{i}(x)\sigma^{j}(x)=0$, and since $D$ is linear, $D(0)=0$. Therefore
$\mathcal M(x,D)=0$ for every $x$ and $D$, and by \noderef{hhessS}, for every $x$, every
$1$-jet $D$ and every $2$-jet $B$,
\[
 \big(\mathrm{hhessS}\ \sigma\ x\ D\ B\big)_{i,j}
  =B\big(\sigma^i(x),\sigma^j(x)\big)+\mathcal M(x,D)_{i,j}=B(e_i,e_j),
\]
again independently of $x$.

\medskip\noindent\textbf{Step 4: the conditions of \noderef{CarnotSetting}.}

\emph{(i) Each $\sigma^i$ is $C^\infty$.} Each $\sigma^i$ is constant, hence $C^\infty$.

\emph{(ii) $\sigma^1(x),\ldots,\sigma^N(x)$ are linearly independent.} They are
$e_1,\ldots,e_N$, the standard basis of $\mathbb R^N$, which is linearly independent.

\emph{(iii) H\"ormander's rank condition (1.6).} Fix $x\in\mathbb R^N$; we must show
that the span of $\{V(x):V\in\mathrm{LieHull}\ \sigma\}$ is all of $\mathbb R^N$. By
\noderef{LieHull}, $\mathrm{LieHull}\ \sigma$ is the intersection of all families of vector
fields that contain each $\sigma^i$ and are closed under the Lie bracket. Every such family
contains the vector field $y\mapsto\sigma^i(y)=e_i$, by the first clause of admissibility;
hence that vector field belongs to the intersection, i.e.
$(y\mapsto e_i)\in\mathrm{LieHull}\ \sigma$ for each $i$. Evaluating at $x$ gives
\[
 \{e_1,\ldots,e_N\}\ \subseteq\ \{V(x):V\in\mathrm{LieHull}\ \sigma\},
\]
and therefore
$\mathbb R^N=\mathrm{span}\{e_1,\ldots,e_N\}\subseteq
\mathrm{span}\{V(x):V\in\mathrm{LieHull}\ \sigma\}\subseteq\mathbb R^N$,
so the two spans coincide with $\mathbb R^N$. Note that only the inclusion
$\sigma^i\in\mathrm{LieHull}\ \sigma$ is used; the hull is never computed. (It is in fact
exhausted at once, since the bracket $[V,W](x)=DW(x)V(x)-DV(x)W(x)$ of two constant fields
is $0$, but this is not needed.)

\emph{(iv) The Taylor formula (2.18).} Let $u:\mathbb R^N\to\mathbb R$ be
differentiable at $x$ with differential $D$. We must show that $u$ is Pansu differentiable
at $x$ with horizontal differential $\mathrm{hgrad}\ \sigma\ x\ D$ in the sense of
\noderef{HasPansuGradAt}: for every $\varepsilon>0$ there is $\delta>0$ such that
\[
 \Big|\,u(x*g)-u(x)-\big\langle \mathrm{hgrad}\ \sigma\ x\ D,\ \pi_1(g)\big\rangle\,\Big|
  \ \le\ \varepsilon\,\|g\|\qquad\text{whenever }\|g\|<\delta .
\]
Here $x*g=x+g$, $\pi_1(g)=g$ and $\|g\|=|g|$. The pairing simplifies: writing
$g=\sum_{i=1}^{N}g_i\,e_i$ and using (3a) and the linearity of $D$,
\[
 \big\langle \mathrm{hgrad}\ \sigma\ x\ D,\ \pi_1(g)\big\rangle
  =\sum_{i=1}^{N}\big(\mathrm{hgrad}\ \sigma\ x\ D\big)_i\,g_i
  =\sum_{i=1}^{N}D(e_i)\,g_i
  =\sum_{i=1}^{N}D\big(g_i\,e_i\big)
  =D\Big(\sum_{i=1}^{N}g_i\,e_i\Big)=D(g).
\]
So the assertion to be proved is: for every $\varepsilon>0$ there is $\delta>0$ with
$|u(x+g)-u(x)-D(g)|\le\varepsilon|g|$ whenever $|g|<\delta$. This is exactly
differentiability of $u$ at $x$ with differential $D$: by definition of the Fr\'echet
derivative, $u(y)-u(x)-D(y-x)=o(|y-x|)$ as $y\to x$, so given $\varepsilon>0$ there is
$\delta>0$ with $|u(y)-u(x)-D(y-x)|\le\varepsilon|y-x|$ for $|y-x|<\delta$; substituting
$y=x+g$ and using $(x+g)-x=g$ gives the claim with the same $\delta$.

\emph{(v) Left-invariance of the horizontal derivatives.} Let $h,x\in\mathbb R^N$, let
$\varphi:\mathbb R^N\to\mathbb R$ and let $(D,B)$ be a second-order jet of $\varphi$ at
$h*x=h+x$ in the sense of \noderef{TwiceDiffAt}. We claim that $(D',B'):=(D,B)$ has the
three required properties for $\varphi\circ L_h$, $L_h(y)=h+y$. (This is the expected
transport $D'=D\circ DL_h(x)$, $B'(u,v)=B(DL_h(x)u,DL_h(x)v)+D(D^2L_h(x)(u,v))$ specialized
to $DL_h(x)=\mathrm{id}$ and $D^2L_h(x)=0$, which is what makes the Euclidean case
immediate.)

\quad\emph{$(D,B)$ is a second-order jet of $\varphi\circ L_h$ at $x$.} The symmetry of $B$
is part of the hypothesis. For the expansion, let $\varepsilon>0$ and take $\delta>0$ from
the hypothesis at $h+x$, so that
\[
 \Big|\varphi(y')-\varphi(h+x)-D\big(y'-(h+x)\big)
   -\tfrac12B\big(y'-(h+x),y'-(h+x)\big)\Big|\le\varepsilon\,\big|y'-(h+x)\big|^{2}
\]
for every $y'$ with $|y'-(h+x)|<\delta$. Now let $y\in\mathbb R^N$ with $|y-x|<\delta$ and
put $y':=h+y$. Then $y'-(h+x)=(h+y)-(h+x)=y-x$, so $|y'-(h+x)|=|y-x|<\delta$, the displayed
inequality applies, and after substituting $y'-(h+x)=y-x$ it reads
\[
 \Big|(\varphi\circ L_h)(y)-(\varphi\circ L_h)(x)-D(y-x)-\tfrac12B(y-x,y-x)\Big|
   \le\varepsilon\,|y-x|^{2},
\]
because $(\varphi\circ L_h)(y)=\varphi(h+y)$ and $(\varphi\circ L_h)(x)=\varphi(h+x)$. This
is the required expansion, with the same $\delta$.

\quad\emph{Equality of the horizontal gradients.} By (3a), the $i$-th component of both
$\mathrm{hgrad}\ \sigma\ x\ D$ and $\mathrm{hgrad}\ \sigma\ (h+x)\ D$ is $D(e_i)$; hence the
two vectors are equal.

\quad\emph{Equality of the symmetrized horizontal second derivatives.} By (3b), the
$(i,j)$ entry of both $\mathrm{hhessS}\ \sigma\ x\ D\ B$ and
$\mathrm{hhessS}\ \sigma\ (h+x)\ D\ B$ is $B(e_i,e_j)$; hence the two matrices are equal.

\emph{(vi) The form of the $X_i$ on functions of the first layer.} Let
$\psi:\mathbb R^{m}\to\mathbb R$ and let $(E,C)$ be a second-order jet of $\psi$ at
$\pi_1(x)$ in the sense of \noderef{TwiceDiffAt}. Since $m=n=N$ and $\pi_1$ is the identity
we have $\pi_1(x)=x$ and $\psi\circ\pi_1=\psi$. Take $(D,B):=(E,C)$. Then
$(D,B)$ is a second-order jet of $\psi\circ\pi_1=\psi$ at $x=\pi_1(x)$, which is the
hypothesis itself. By (3a) the $i$-th component of $\mathrm{hgrad}\ \sigma\ x\ E$ is
$E(e_i)$, so $\mathrm{hgrad}\ \sigma\ x\ E=(E(e_i))_{i=1}^{m}$ as required; and by (3b) the
$(i,j)$ entry of $\mathrm{hhessS}\ \sigma\ x\ E\ C$ is $C(e_i,e_j)$, so
$\mathrm{hhessS}\ \sigma\ x\ E\ C=(C(e_i,e_j))_{i,j=1}^{m}$ as required.

\medskip\noindent\textbf{Step 5: the five listed properties.} By construction $n=N$,
$m=N$ and $r=1$, which is (1); $x*y=x+y$, which is (2); and $\|x\|=|x|$, which is (3). For
(4), $\exp_{\mathrm h}$ and $\pi_1$ are both the identity, so
$\exp_{\mathrm h}(\pi_1(x))=x$. For (5), $\sigma^i(x)=e_i$ and
$\exp_{\mathrm h}(e_i)=e_i$, so $\sigma^i(x)=\exp_{\mathrm h}(e_i)$.

This exhibits the required Carnot setting and completes the proof.

\medskip\noindent\emph{Remark on the role of this node.} Item 3 of the paper's note for the
formalization permits the structural facts of Section Section~2.4 to be built into the
definition of the setting only ``provided every Carnot group in the coordinates of Section
Section~2.4 satisfies that definition''. The theorem above discharges that proviso at the
one instance that can be checked outright, and it is the instance Theorem Theorem~1.3(iii)
needs: the comparison principle for $r=1$ is applied in $\mathbb R^{m}$, on the
abelianization, so a Carnot setting on $\mathbb R^{m}$ with $x*y=x+y$ and
$X_i=\partial_{x_i}$ must exist for that step to be available at all. Properties (4) and
(5) record how that setting is identified with the Euclidean one: by (4) and the field
$\pi_1(\exp_{\mathrm h}(\zeta))=\zeta$ of \noderef{CarnotGroup}, the continuous linear map
$\pi_1$ is a bijection with two-sided inverse $\exp_{\mathrm h}$, so $\exp_{\mathrm h}$ is
linear as well; by (3) and the field $\|\exp_{\mathrm h}(\zeta)\|=|\zeta|$ of
\noderef{CarnotNorm} it is an isometry; and by (5) the frame is the image
$\exp_{\mathrm h}(e_1),\ldots,\exp_{\mathrm h}(e_m)$ of the standard basis, so by (3a) and
(3b) the horizontal gradient and the symmetrized horizontal second derivative of a jet are
its Euclidean gradient and Hessian read in that orthonormal basis. For $N=0$ the setting is
the trivial group, which is a Carnot group of step $1$ with no generators; the case used in
the paper is $N\ge1$.
\end{proof}
